\documentclass[exercice,a4]{thedocument}%
\usepackage{texmaths-tikz}%
\usepackage{tabularx}%
\usepackage[np]{numprint}%
\startdatakeys
\source{24GENMATG11}
\BO{2026}
\cycle{4}
\competencesDuSocle{CA,RA}
\objectifsApprentissage{C4B2cb,C4B3cb,C4D2ac,C4D2ae,C4D2af}
\motsCles{cône,Pythagore,patron,arc,cercle,volume,pourcentage}
\enddatakeys
\begin{document}%
\def\wideparen#1{\overset{\frown}{#1}}%
\textbf{Partie A}

Léo veut fabriquer un chapeau en forme de cône. Le rayon OM de la base
mesure $\np[cm]{9}$ et la hauteur OS mesure $\np[cm]{30}$.

\begin{center}
\begin{tikzpicture}
    \draw(0,0)ellipse(2 and .5);
    \draw(2,0)--(0,5)--(-2,0);
    \draw[dashed](0,0)--(0,5);
    \node[left]at(0,0){O};\node[right]at(2,0){M};\node[above]at(0,5){S};
\end{tikzpicture}
\end{center}

\begin{enumerate}
    \item Démontrer que la longueur MS, arrondie au dixième de centimètre,
    est $\np[cm]{31,3}$.
    \item Léo a un tour de tête de $\np[cm]{56}$. Les dimensions choisies
    sont-elles adaptées ?
    \item Le patron est un secteur de cercle de rayon $SM$. L'arc
    $\wideparen{\mathrm{M'M}}$ mesure $\np[cm]{56,5}$.
    \begin{enumerate}
        \item Démontrer que le cercle de centre S et de rayon SM a une
        longueur, arrondie au dixième, de $\np[cm]{196,7}$.
        \item Calculer la mesure de l'angle $\widehat{\mathrm{M'SM}}$
        correspondant à cet arc, arrondie au degré.
    \end{enumerate}
\end{enumerate}

\textbf{Partie B}

On rappelle que la hauteur du chapeau mesure $\np[cm]{30}$.
\begin{enumerate}
    \item Montrer que le volume total du chapeau, arrondi au $\mathrm{cm^3}$,
    est $\np[cm^3]{2545}$.

    \begin{center}
    \fbox{\begin{minipage}{.8\linewidth}
    Le volume d'un cône de rayon $R$ et de hauteur $h$ est
    \[
        V=\dfrac13\times(\pi R^2)\times h.
    \]
    \end{minipage}}
    \end{center}

    \item Les bonbons atteignent le milieu de la hauteur du chapeau.
    Léo pense que leur volume représente moins de 15\,\% du volume total.
    Est-ce exact ? Justifier.
\end{enumerate}
\showthesource
\end{document}
